\chapter{Języki formalne, termy i formuły}\label{ch:syntax-fol}

\section{Dlaczego odróżniamy napis od jego znaczenia?}
Zdanie „dla każdej liczby istnieje większa” jest zrozumiałe, lecz pozostawia wiele
rzeczy domyślnych: z jakiego zbioru pochodzą liczby, co znaczy „większa” i jaki
jest zakres kwantyfikatorów. W dowodzie zwykle radzimy sobie z tymi kwestiami,
ale gdy przedmiotem badań staje się \emph{sam dowód}, potrzebujemy języka,
którego poprawne wyrażenia można rozpoznawać według jednoznacznych reguł.
Pytanie, czy wyrażenie zostało dobrze zapisane, należy do \emph{składni};
pytanie, czy jest prawdziwe, należy do \emph{semantyki}.

Warto pamiętać o dwóch poziomach wypowiedzi. Pisząc „$\varphi$ jest formułą”,
mówimy \emph{o} pewnym napisie, czyli używamy metajęzyka. Pisząc symbol $\forall$
wewnątrz $\varphi$, tworzymy wyrażenie samego języka formalnego. W dalszych
rozdziałach rozróżnienie to okaże się niezbędne do sformułowania twierdzeń Gödla.

\section{Alfabety, słowa i sygnatury}
\begin{definicja}[Alfabet i słowo]\label{def:syntax-alphabet}
\emph{Alfabetem} nazywamy zbiór symboli $\Sigma$. \emph{Słowem} nad $\Sigma$
nazywamy skończony ciąg $(s_1,\ldots,s_n)$ symboli $s_i\in\Sigma$.
Dopuszczamy $n=0$: otrzymujemy wtedy \emph{słowo puste}, oznaczane $\varepsilon$.
Zbiór wszystkich słów nad $\Sigma$ oznaczamy przez $\Sigma^*$.
\emph{Konkatenacją} słów $u$ i $v$ jest słowo $uv$, otrzymane przez dopisanie
symboli $v$ za symbolami $u$.
\end{definicja}
\begin{przyklad}
Dla $\Sigma=\{a,b\}$ słowa $ab$, $aab$ i $\varepsilon$ należą do $\Sigma^*$.
Mamy $a\varepsilon=\varepsilon a=a$, natomiast $ab\ne ba$.
Słowo samo przez się nie ma jeszcze znaczenia matematycznego.
\end{przyklad}
\begin{stwierdzenie}[Łączność konkatenacji]\label{prop:syntax-concat}
Niech $\Sigma$ będzie alfabetem. Dla dowolnych słów $u,v,w\in\Sigma^*$ zachodzi $(uv)w=u(vw)$.
\end{stwierdzenie}
\begin{proof}
Niech $u=(u_1,\ldots,u_k)$, $v=(v_1,\ldots,v_\ell)$ i
$w=(w_1,\ldots,w_m)$. W obu sposobach łączenia otrzymujemy ciąg długości
$k+\ell+m$, którego pierwsze $k$ pozycji zajmują symbole $u_i$, następne
$\ell$ pozycji symbole $v_j$, a ostatnie $m$ pozycji symbole $w_r$.
Dwa skończone ciągi o jednakowej długości i jednakowych symbolach na każdej
pozycji są równe, co kończy dowód.
\end{proof}

\begin{twierdzenie}[Słowa tworzą monoid wolny]\label{thm:syntax-free-monoid-xlv}
Dla dowolnego alfabetu $\Sigma$ konkatenacja jest działaniem
łącznym na $\Sigma^*$, a słowo puste $\varepsilon$ jest jej
obustronnym elementem neutralnym. Ponadto funkcja długości
$|\cdot|:\Sigma^*\to\N$ spełnia $|uv|=|u|+|v|$.
Dla każdego monoidu $(M,\cdot,e)$ i każdej funkcji
$j:\Sigma\to M$ istnieje dokładnie jeden homomorfizm monoidów
$\overline j:\Sigma^*\to M$ taki, że
$\overline j((a))=j(a)$ dla każdej litery $a\in\Sigma$.
\end{twierdzenie}
\begin{proof}
Łączność wykazano w Stwierdzeniu~\ref{prop:syntax-concat}.
Jeżeli $u=(u_1,\ldots,u_k)$, to dopisanie na końcu zera symboli
nie zmienia żadnej spośród $k$ pozycji, stąd $u\varepsilon=u$.
Dopisanie zera symboli przed pierwszą pozycją daje
$\varepsilon u=u$; twierdzenie obejmuje też $k=0$.
Z definicji konkatenacji $uv$ zawiera najpierw $|u|$ symboli,
a następnie $|v|$ symboli. Te dwa kolejne bloki nie
nakładają się, więc łączna długość jest ich sumą.

Niech teraz $(M,\cdot,e)$ będzie dowolnym monoidem, a
$j:\Sigma\to M$ funkcją. Definiujemy
$\overline j(\varepsilon)=e$ oraz
$\overline j((a_1,\ldots,a_n))=
 j(a_1)\cdot\ldots\cdot j(a_n)$ dla $n\geq1$.
Łączność działania w $M$ zapewnia niezależność wyniku
od ustawienia nawiasów. Dla słów $u,v$ ich konkatenacja
polega na dopisaniu listy liter $v$ do listy liter $u$,
więc przez łączność i własność elementu neutralnego
$\overline j(uv)=\overline j(u)\cdot\overline j(v)$,
także gdy jedno ze słów jest puste.
Otrzymaliśmy homomorfizm, który na słowach jednoliterowych
zgadza się z $j$.
Jeżeli $h$ jest innym takim homomorfizmem, to
$h(\varepsilon)=e$, a dla każdego niepustego słowa
$w=(a_1)\cdots(a_n)$ mamy
$h(w)=h((a_1))\cdots h((a_n))
=j(a_1)\cdots j(a_n)=\overline j(w)$.
Zatem homomorfizm jest jednoznaczny, co uzasadnia
określenie \emph{monoid wolny na alfabecie $\Sigma$}.
\end{proof}
\begin{przyklad}[Czy kolejność dopisywania ma znaczenie?]
Dla $\Sigma=\{a,b\}$, $u=ab$, $v=b$ i $w=a$ otrzymujemy
$(uv)w=abba=u(vw)$, a $|uv|=3=2+1$.
Nie zachodzi natomiast przemienność: $uv=abb$, podczas
gdy $vu=bab$. Słowo $\varepsilon$ nie jest literą alfabetu,
lecz ciągiem o długości zero.
\end{przyklad}

\begin{definicja}[Sygnatura języka pierwszego rzędu]\label{def:syntax-signature}
\emph{Sygnaturą} $\mathcal L$ nazywamy zbiór symboli pozalogicznych
z podziałem na symbole stałych $\operatorname{Const}(\mathcal L)$,
funkcyjne $\operatorname{Fun}(\mathcal L)$ i relacyjne
$\operatorname{Rel}(\mathcal L)$. Każdemu symbolowi funkcyjnemu i relacyjnemu
przypisujemy dodatnią liczbę naturalną, zwaną \emph{arnością};
symbole stałych mają arność $0$.
Symbole należące do różnych kategorii są odrębne, nawet jeśli w zapisie
nieformalnym używamy dla nich tych samych liter.
\end{definicja}
Wybieramy ponadto nieskończony ciąg zmiennych $v_0,v_1,\ldots$ oraz znaki
$=,\neg,\land,\lor,\to,\leftrightarrow,\forall,\exists$ i symbole pomocnicze
(nawiasy, przecinki). Znak równości traktujemy jako symbol logiczny.
\emph{Językiem pierwszego rzędu o sygnaturze $\mathcal L$} nazywamy te
symbole wraz z opisanymi poniżej regułami tworzenia termów i formuł.

\begin{przyklad}[Trzy sygnatury]\label{ex:syntax-signatures}
W języku arytmetyki mamy stałą $0$, symbol następstwa $S$ arności $1$,
symbole $+$, $\cdot$ arności $2$ i symbol relacyjny $<$ arności $2$.
Język grafów może zawierać tylko relację $E$ arności $2$.
W języku grup używamy stałej $e$, działania $\ast$ arności $2$
i symbolu $\iota$ arności $1$ dla operacji odwracania.
Żaden z tych symboli nie ma jeszcze przypisanej interpretacji.
\end{przyklad}


\begin{uwaga}[Dlaczego pojawia się teraz symbol $\to$?]
W początkowych rozdziałach zapis $p\Rightarrow q$ oznaczał
implikację między zdaniami. W ścisłym języku formalnym używamy
symbolu $\to$ jako \emph{znaku budującego formułę}.
Gdy zaś w zwykłym tekście matematycznym napiszemy
„z $A$ wynika $B$”, jest to wypowiedź o związku między formułami,
a niekoniecznie kolejna formuła wewnątrz języka.
Rozróżnienie symbolu wewnątrz języka i wypowiedzi o języku
będzie szczególnie ważne w twierdzeniu o pełności.
\end{uwaga}

\section{Termy: napisy nazywające obiekty}
\begin{definicja}[Term]\label{def:syntax-term}
Zbiór termów $\operatorname{Term}(\mathcal L)$ jest najmniejszym zbiorem
napisów spełniającym następujące warunki:
\begin{enumerate}[label=(\roman*)]
\item każda zmienna $v_i$ jest termem;
\item każdy symbol stałej jest termem;
\item jeśli $f$ ma arność $n$ i $t_1,\ldots,t_n$ są termami,
      to $f(t_1,\ldots,t_n)$ jest termem.
\end{enumerate}
Innych termów nie ma. \emph{Term domknięty} to term bez wystąpień zmiennych.
W zapisie formalnym symbole działań występują przed argumentami; zapis
$+(t_1,t_2)$ zwyczajowo skracamy do $t_1+t_2$.
\end{definicja}
\begin{przyklad}[Co jest, a co nie jest termem?]
Dla sygnatury arytmetyki termami są $S(0)$, $+(v_2,S(0))$
i $\cdot(+(v_0,0),v_3)$. Napisy $S(v_0,v_1)$ i $+(,v_0)$
nie są termami: pierwszy zawiera za dużo argumentów, drugi nie zawiera
pierwszego argumentu. Napis $v_0<v_1$ nie jest termem, lecz formułą atomową.
\end{przyklad}

\begin{lemat}[Jednoznaczne odczytywanie termów]\label{lem:syntax-unique-terms}
Niech $\mathcal L$ będzie sygnaturą. Każdy term języka $\mathcal L$ jest dokładnie jednego z trzech rodzajów:
zmienną, stałą lub napisem $f(t_1,\ldots,t_n)$.
W trzecim przypadku symbol $f$ i wszystkie termy $t_1,\ldots,t_n$
są jednoznacznie wyznaczone.
\end{lemat}
\begin{proof}
Trzy rodzaje rozpoczynają się różnymi rodzajami symboli, więc nie można
pomieszać zmiennej, stałej i zastosowania symbolu funkcyjnego.
Jeśli jeden napis da się zapisać jako $f(t_1,\ldots,t_n)$ i
$g(s_1,\ldots,s_m)$, to porównanie jego pierwszych symboli daje $f=g$,
a przypisana arność wymusza $n=m$.
Przecinki rozdzielające argumenty znajdują się na poziomie pierwszej pary
nawiasów po $f$. Przecinków należących do termów zagnieżdżonych nie należy
brać pod uwagę: występują one przy dodatniej głębokości zagnieżdżenia.
Zatem miejsce zakończenia każdego argumentu jest jednoznaczne.
Porównując odpowiadające im fragmenty obu identycznych napisów, otrzymujemy
$t_i=s_i$ dla każdego $i$. Formalne uzasadnienie można prowadzić
indukcyjnie po długości termu.
\end{proof}

\begin{przyklad}[Jednoznaczny rozkład termu]
W sygnaturze z dwuargumentowym $f$, jednoargumentowym $g$
i stałą $c$ rozważmy $t=f(g(v_0),f(c,v_1))$.
Zewnętrzny symbol $f$ ma dwa argumenty. Pierwszy to
$g(v_0)$, drugi to $f(c,v_1)$; przecinek wewnątrz
drugiego argumentu znajduje się na wyższym poziomie
zagnieżdżenia i nie rozdziela argumentów zewnętrznego $f$.
Drzewo konstrukcji ma liście $v_0,c,v_1$, a jego korzeń
ma etykietę $f$. Jest więc jeden sposób odczytania termu
zgodny z arnościami ustalonej sygnatury.
\end{przyklad}

\section{Formuły: napisy wyrażające własności}
\begin{definicja}[Formuła atomowa i formuła]\label{def:syntax-formula}
\emph{Formułą atomową} jest napis $t_1=t_2$ albo $R(t_1,\ldots,t_n)$,
gdzie $t_i$ są termami, a $R$ symbolem relacyjnym arności $n$.
Zbiór formuł jest najmniejszym zbiorem zawierającym wszystkie formuły atomowe
i zamkniętym na operacje
\[
\varphi\longmapsto\neg\varphi,
\qquad
(\varphi,\psi)\longmapsto
(\varphi\circ\psi),\quad
\circ\in\{\land,\lor,\to,\leftrightarrow\},
\]
oraz na tworzenie $\forall v_i\,\varphi$ i $\exists v_i\,\varphi$
dla dowolnego $i\in\N$. Nic innego nie jest formułą.
Dla jednoznaczności parsowania formuły z głównym spójnikiem binarnym
w pełnym zapisie są objęte nawiasami. W przykładach używamy również
nawiasów grupujących dla czytelności, na przykład zapis
$\forall v_0\,(v_0=v_0)$ oznacza ściśle tę samą formułę
co $\forall v_0\,v_0=v_0$. Nadmiarowe nawiasy grupujące
nie stanowią dodatkowej reguły tworzenia formuł i nie są
częścią napisu w ścisłej gramatyce. Tak rozumiane skróty
nie wpływają na jednoznaczność jego odczytywania.
\end{definicja}
\begin{przyklad}
W języku arytmetyki napis $\forall v_0\,(v_0=v_0)$ jest zdaniem
w sensie zdefiniowanym poniżej, natomiast
$\exists v_1\,(+(v_1,v_1)=v_0)$ jest formułą z wolną zmienną $v_0$.
Napisy $\forall(v_0=v_0)$ oraz $\exists v_1\,+(v_1,v_1)$ nie są formułami.
\end{przyklad}

\begin{lemat}[Jednoznaczność budowy formuły]\label{lem:syntax-unique-formulas}
Niech $\mathcal L$ będzie sygnaturą. Dla każdej formuły języka $\mathcal L$ $\varphi$ można jednoznacznie odtworzyć jej ostatni
krok konstrukcji, a następnie wszystkie bezpośrednie podformuły.
\end{lemat}
\begin{proof}
Jeśli napis zaczyna się od symbolu relacyjnego, jego arność i nawiasy
określają argumenty. Jeśli zaczyna się od termu, to po jego końcu
w formule atomowej musi wystąpić znak równości i drugi term; jednoznaczność
termów zapewnia Lemat~\ref{lem:syntax-unique-terms}.
Jeśli napis zaczyna się od $\neg$, usunięcie pierwszego symbolu daje
jedyną podformułę. Jeśli zaczyna się od kwantyfikatora, następny symbol
wskazuje związaną zmienną, a pozostały napis jest podformułą.
Pozostają nawiasowane formuły binarne. Czytając je od lewej do prawej
i zliczając poziom otwarcia nawiasów, znajdujemy dokładnie jeden
spójnik na zewnętrznym poziomie. Dzieli on napis na lewą i prawą
podformułę. Nie może być ich dwóch: wewnętrzne spójniki znajdują
się w zagnieżdżonych nawiasach, a formuła musi być zbudowana według
podanej gramatyki. Powtarzanie procedury dla podformuł kończy się,
ponieważ są one krótsze od formuły wyjściowej.
\end{proof}

\begin{przyklad}[Drzewo budowy formuły]\label{ex:syntax-formula-tree-xlv}
Dla relacji dwuargumentowej $R$ i symbolu funkcyjnego
jednoargumentowego $f$ rozważmy
$\theta=\forall v_0\,(R(v_0,f(v_1))\to v_0=v_1)$.
Ostatnią operacją jest kwantyfikacja po $v_0$; pod nią
znajduje się implikacja z dwiema podformułami atomowymi.
Drzewo odzwierciedla kolejność konstrukcji, nie jest
rysunkiem wartościowania ani modelem relacji $R$.
\end{przyklad}
\begin{figure}[htbp]
\centering
\begin{tikzpicture}[>=Latex, every node/.style={draw,rounded corners=2pt,inner sep=4pt,font=\small}]
  \node (q) at (0,1.9) {$\forall v_0$};
  \node (imp) at (0,0.95) {$\to$};
  \node (rel) at (-2.45,0) {$R(v_0,f(v_1))$};
  \node (eq) at (2.15,0) {$v_0=v_1$};
  \draw[->] (q)--(imp);
  \draw[->] (imp)--(rel);
  \draw[->] (imp)--(eq);
\end{tikzpicture}
\caption{Drzewo składniowe formuły $\theta$. Strzałki prowadzą
od konstrukcji do jej bezpośrednich podformuł.}
\label{fig:syntax-tree-xlv}
\end{figure}

\section{Zmienne wolne, związane i zakres kwantyfikatora}
\begin{definicja}[Zmienne wolne]\label{def:syntax-free-vars}
Zbiór $\operatorname{FV}(t)$ zmiennych termu definiujemy przez
$\operatorname{FV}(v_i)=\{v_i\}$, $\operatorname{FV}(c)=\varnothing$
i
$\operatorname{FV}(f(t_1,\ldots,t_n))=\bigcup_{j=1}^{n}\operatorname{FV}(t_j)$.
Dla formuł definiujemy
\begin{align*}
\operatorname{FV}(t_1=t_2)
  &=\operatorname{FV}(t_1)\cup\operatorname{FV}(t_2),\\
\operatorname{FV}(R(t_1,\ldots,t_n))
  &=\bigcup_{j=1}^{n}\operatorname{FV}(t_j),\\
\operatorname{FV}(\neg\varphi)&=\operatorname{FV}(\varphi),\\
\operatorname{FV}(\varphi\circ\psi)
  &=\operatorname{FV}(\varphi)\cup\operatorname{FV}(\psi),\\
\operatorname{FV}(\forall v_i\,\varphi)
  &=\operatorname{FV}(\varphi)\setminus\{v_i\},\\
\operatorname{FV}(\exists v_i\,\varphi)
  &=\operatorname{FV}(\varphi)\setminus\{v_i\}.
\end{align*}
W tych wzorach $\circ$ oznacza dowolny spójnik binarny.
Formułę bez zmiennych wolnych nazywamy \emph{zdaniem}.
\end{definicja}
\begin{przyklad}[To samo oznaczenie może wystąpić wolno i związanie]
W formule $(\exists v_1\,(v_0=v_1))\land(v_1=v_2)$
pierwsze wystąpienie $v_1$ po kwantyfikatorze jest związane, ale
w drugim członie $v_1$ występuje wolno. Zatem
$\operatorname{FV}(\varphi)=\{v_0,v_1,v_2\}$.
Kwantyfikator nie oddziałuje na wyrażenie znajdujące się poza
zakresem swojej podformuły.
\end{przyklad}

\begin{przyklad}[Obliczanie zmiennych wolnych etapami]
Niech
$\rho=(\exists v_1\,R(v_0,v_1))\land
 (\forall v_0\,(R(v_0,v_2)\lor R(v_1,v_0)))$,
gdzie $R$ jest symbolem relacyjnym arności $2$.
Pierwszy człon ma zmienne wolne $\{v_0\}$: $v_1$ zostało
związane. Podformuły atomowe drugiego członu mają
zbiory zmiennych odpowiednio $\{v_0,v_2\}$ i
$\{v_1,v_0\}$, więc ich alternatywa ma zbiór
$\{v_0,v_1,v_2\}$. Kwantyfikator $\forall v_0$
usuwa z niego $v_0$ i pozostawia $\{v_1,v_2\}$.
Ostatecznie $\operatorname{FV}(\rho)=\{v_0,v_1,v_2\}$.
Wystąpienie zmiennej $v_1$ jest związane w pierwszym
członie, lecz wolne w drugim; podobnie $v_0$ jest wolne
w pierwszym członie, ale związane w drugim.
\end{przyklad}

\section{Podstawianie bez przechwytywania zmiennych}
Chcemy wstawić do formuły term opisujący konkretny obiekt.
Nie wystarczy jednak zastąpić każdego wystąpienia danej litery:
kwantyfikator może przez przypadek związać zmienną podstawianego termu.

\begin{definicja}[Podstawienie termu w termie]\label{def:syntax-term-substitution-xlv}
Dla termów $u,t$ języka $\mathcal L$ i zmiennej $x$
term $t[x:=u]$ definiujemy rekurencyjnie:
\begin{enumerate}[label=(\roman*)]
\item $x[x:=u]=u$, a $y[x:=u]=y$ dla każdej zmiennej $y\ne x$;
\item $c[x:=u]=c$ dla każdego symbolu stałej $c$;
\item $f(t_1,\ldots,t_n)[x:=u]
 =f(t_1[x:=u],\ldots,t_n[x:=u])$.
\end{enumerate}
Termy nie zawierają kwantyfikatorów, więc żadna zmienna
podstawianego termu nie może zostać w ten sposób związana.
\end{definicja}
\begin{lemat}[Zmienne termu po podstawieniu]
\label{lem:syntax-term-fv-subst-xlv}
Niech $\mathcal L$ będzie sygnaturą. Dla dowolnych termów
$t,u$ języka $\mathcal L$ i zmiennej $x$ zachodzi
\[
\operatorname{FV}(t[x:=u])=
\begin{cases}
(\operatorname{FV}(t)\setminus\{x\})\cup\operatorname{FV}(u),
 &x\in\operatorname{FV}(t),\\
\operatorname{FV}(t),&x\notin\operatorname{FV}(t).
\end{cases}
\]
\end{lemat}
\begin{proof}
W przypadku $t=x$ wynikiem jest $u$, a po prawej
otrzymujemy $\operatorname{FV}(u)$, więc równość zachodzi.
Dla $t=y\ne x$ nie zmieniamy termu ani jego jedynej
zmiennej. Dla stałej oba zbiory zmiennych są puste.
Przy $t=f(t_1,\ldots,t_n)$ podstawienie wykonuje się
oddzielnie w każdym argumencie. Jeżeli $x$ nie występuje
w żadnym $t_i$, nie zmienia się żaden argument.
W przeciwnym razie dla argumentów zawierających $x$
założenie indukcyjne usuwa tę zmienną i dodaje cały zbiór
$\operatorname{FV}(u)$, a pozostałe argumenty zachowują swoje
zmienne. Po zsumowaniu zbiorów zmiennych wszystkich
argumentów otrzymujemy dokładnie pierwszy przypadek wzoru.
\end{proof}
\begin{przyklad}[Podstawienie w zagnieżdżonym termie]
Niech $t=f(x,g(x,y))$ i $u=h(y)$ dla symboli funkcyjnych
$f$ arności $2$, $g$ arności $2$ i $h$ arności $1$.
Wtedy $t[x:=u]=f(h(y),g(h(y),y))$.
Obie początkowe kopie $x$ są zastąpione; nowy term ma
zbiór zmiennych $\{y\}$, zgodnie z lematem.
\end{przyklad}

\begin{definicja}[Dopuszczalne podstawienie]\label{def:syntax-free-for}
Niech $x$ będzie zmienną, $t$ termem, a $\varphi$ formułą.
Mówimy, że $t$ jest \emph{dopuszczalny do podstawienia za $x$ w $\varphi$},
jeśli żadne wolne wystąpienie $x$ nie znajduje się w zakresie
kwantyfikatora wiążącego zmienną z $\operatorname{FV}(t)$.
Wtedy $\varphi[x:=t]$ oznacza formułę powstałą przez zastąpienie
\emph{wolnych} wystąpień $x$ termem $t$. Formalnie:
na atomach podstawiamy term w każdym jego argumencie,
a dla negacji i spójników binarnych wykonujemy operację
w bezpośrednich podformułach. Przy $Qy\,\psi$, gdzie
$Q\in\{\forall,\exists\}$, jeżeli $y=x$, pozostawiamy
całą formułę bez zmian. Jeżeli $y\ne x$, podstawiamy
rekurencyjnie w $\psi$ i zachowujemy kwantyfikator $Qy$.
Ta ostatnia reguła jest dozwolona tylko przy spełnieniu
powyższego warunku braku przechwycenia; w przeciwnym razie
najpierw trzeba poprawnie przemianować odpowiednią
zmienną związaną. Jeżeli $x\notin\operatorname{FV}(\varphi)$,
podstawienie nie zmienia formuły.
\end{definicja}
\begin{przyklad}[Przechwycenie zmiennej]
Rozważmy $\varphi=\forall v_1\,(v_0=v_1)$ i term $t=v_1$.
Bezpośrednia zamiana $v_0$ na $v_1$ dałaby
$\forall v_1\,(v_1=v_1)$, w którym nowe wystąpienie $v_1$
stało się związane. Podstawienie nie jest dopuszczalne.
Możemy jednak najpierw zmienić nazwę zmiennej związanej na świeżą
zmienną $v_2$, otrzymując $\forall v_2\,(v_0=v_2)$, a następnie
wykonać dopuszczalne podstawienie i otrzymać
$\forall v_2\,(v_1=v_2)$.
\end{przyklad}
\begin{lemat}[Zmienne wolne po podstawieniu]\label{lem:syntax-fv-subst}
Niech $\mathcal L$ będzie sygnaturą, $x$ zmienną,
$t$ termem języka $\mathcal L$, a $\varphi$ formułą tego języka.
Załóżmy, że $t$ jest dopuszczalny do podstawienia za $x$ w $\varphi$.
Jeśli $x\in\operatorname{FV}(\varphi)$, to
\[
\operatorname{FV}(\varphi[x:=t])
=(\operatorname{FV}(\varphi)\setminus\{x\})
\cup\operatorname{FV}(t).
\]
Jeśli $x\notin\operatorname{FV}(\varphi)$, podstawienie nie zmienia
formuły ani zbioru jej zmiennych wolnych.
\end{lemat}
\begin{proof}
Dowodzimy tezy indukcją po budowie $\varphi$.
Dla atomów równościowych i relacyjnych operacja dotyczy
wyłącznie termów. Wzór wynika wtedy z
Lematu~\ref{lem:syntax-term-fv-subst-xlv} po wzięciu sumy
zbiorów zmiennych argumentów. W przypadku negacji
$\operatorname{FV}(\neg\psi)=\operatorname{FV}(\psi)$,
więc teza jest dokładnie założeniem indukcyjnym.
Dla spójnika binarnego stosujemy to założenie osobno
w dwóch podformułach i sumujemy otrzymane zbiory.
Jeżeli w całej formule $x$ nie jest wolna, definicja
podstawienia nie zmienia żadnego wolnego wystąpienia $x$;
formuła pozostaje ta sama.

Rozważmy przypadek $\varphi=Qy\,\psi$, gdzie
$Q\in\{\forall,\exists\}$ i $x$ jest wolna w $\varphi$.
Wówczas $x\ne y$ oraz $x\in\operatorname{FV}(\psi)$.
Dopuszczalność podstawienia daje
$y\notin\operatorname{FV}(t)$: każde wolne wystąpienie $x$
w podformule $\psi$ jest bowiem w zakresie kwantyfikatora
$Qy$. Zatem, korzystając kolejno z definicji wolnych
zmiennych, założenia indukcyjnego i $y\notin\operatorname{FV}(t)$,
otrzymujemy
\begin{align*}
\operatorname{FV}((Qy\,\psi)[x:=t])
 &=\operatorname{FV}(\psi[x:=t])\setminus\{y\}\\
 &=((\operatorname{FV}(\psi)\setminus\{x\})
    \cup\operatorname{FV}(t))\setminus\{y\}\\
 &=(\operatorname{FV}(\varphi)\setminus\{x\})
    \cup\operatorname{FV}(t).
\end{align*}
To zamyka również oba przypadki kwantyfikatorów.
\end{proof}

\begin{uwaga}[Nazwy zmiennych związanych]
Zmiana nazwy zmiennej związanej na nową, niewystępującą
w rozważanej formule, nie powinna zmieniać jej znaczenia.
Takie formuły nazywa się \emph{$\alpha$-równoważnymi}.
W następnym rozdziale uzasadnimy to na podstawie definicji spełniania,
a w rozdziale o rachunku dowodowym będziemy korzystać wyłącznie
z podstawień dopuszczalnych albo wcześniej zmienimy nazwę zmiennej związanej.
\end{uwaga}

\begin{uwaga}[Dlaczego warunek dopuszczalności jest konieczny?]
Formuła $\forall y\,(x=y)$ ma wolną zmienną $x$.
Próba wstawienia termu $y$ bez przemianowania zmieniłaby
ją w $\forall y\,(y=y)$, czyli w zdanie. Pierwsza formuła
jest fałszywa przy dowolnym wartościowaniu w strukturze
co najmniej dwuelementowej; druga jest w tej strukturze
prawdziwa. Problem nie jest więc tylko notacyjny.
\end{uwaga}

\begin{lemat}[Podstawienie i kwantyfikator zasłaniający]\label{lem:syntax-shadowing-lxvii}
Niech $x,y$ będą zmiennymi, $t$ termem, a $\psi$ formułą.
Jeżeli $x=y$, to $(Qx\,\psi)[x:=t]=Qx\,\psi$ dla
$Q\in\{\forall,\exists\}$. Jeżeli $x\ne y$ i
$x\notin\operatorname{FV}(\psi)$, to również
$(Qy\,\psi)[x:=t]=Qy\,\psi$, niezależnie od zmiennych termu $t$.
Jeżeli $x\ne y$ i $x\in\operatorname{FV}(\psi)$, to
$t$ jest dopuszczalny do podstawienia za $x$ w $Qy\,\psi$
wtedy i tylko wtedy, gdy jest dopuszczalny w $\psi$
oraz $y\notin\operatorname{FV}(t)$.
\end{lemat}
\begin{proof}
W pierwszym przypadku zewnętrzny kwantyfikator wiąże
wszystkie wystąpienia $x$ w swoim zakresie, które mogłyby
być wolne w $\psi$; reguła podstawiania pozostawia całą
formułę bez zmian. W drugim przypadku w $\psi$ nie ma
wolnych wystąpień $x$, więc nie ma czego zastępować.
W trzecim przypadku każde wolne wystąpienie $x$ w $\psi$
leży w zakresie nowego kwantyfikatora $Qy$.
Dopuszczalność w $Qy\,\psi$ wymaga zatem, po pierwsze,
aby żaden kwantyfikator wewnątrz $\psi$ nie przechwytywał
zmiennej z $\operatorname{FV}(t)$, czyli dopuszczalności
w $\psi$, a po drugie, aby $y$ nie była wolna w $t$.
Oba warunki razem wykluczają wszystkie możliwe
przechwycenia i są wystarczające.
\end{proof}
\begin{przyklad}[Wystąpienie wolne i zasłonięte]
Niech $\psi=(x=z)\land\forall x\,(x=y)$, a $t=f(y)$.
W formule $\psi$ wolne wystąpienie $x$ znajduje się tylko
w pierwszym członie. Podstawienie daje
$\psi[x:=t]=(f(y)=z)\land\forall x\,(x=y)$;
wystąpienie $x$ pod wewnętrznym kwantyfikatorem pozostaje
nietknięte. Natomiast w $\forall y\,\psi$ to samo
podstawienie nie jest dopuszczalne, ponieważ $y$ z termu
$f(y)$ zostałaby przechwycona przez zewnętrzny kwantyfikator.
W $\forall x\,\psi$ podstawienie nie zmienia niczego.
\end{przyklad}

\section{Indukcja po budowie formuły}
Definicje termów i formuł podają również sposób prowadzenia dowodów.
Aby wykazać własność wszystkich formuł, wystarczy udowodnić ją
najpierw dla formuł atomowych, a potem sprawdzić, że przechodzi
przez każdy spójnik i oba kwantyfikatory. Nie jest to nowe
założenie logiczne, lecz zastosowanie zasady indukcji do skończonych
drzew konstrukcji napisów.

\begin{twierdzenie}[Zasada indukcji strukturalnej]
\label{thm:syntax-structural-induction-xlv}
Niech $P$ będzie własnością termów języka $\mathcal L$.
Jeśli $P$ zachodzi dla każdej zmiennej i każdej stałej,
a z $P(t_1),\ldots,P(t_n)$ wynika $P(f(t_1,\ldots,t_n))$
dla każdego symbolu funkcyjnego $f$ arności $n$, to
$P$ zachodzi dla wszystkich termów.
Analogicznie niech $Q$ będzie własnością formuł.
Jeśli $Q$ zachodzi dla wszystkich atomów oraz zachowuje się
przy negacji, czterech spójnikach binarnych i obu
kwantyfikatorach, to $Q$ zachodzi dla wszystkich formuł.
\end{twierdzenie}
\begin{proof}
Każda formuła i każdy term powstają w skończenie wielu
krokach, ponieważ zbiory z Definicji~\ref{def:syntax-term}
i Definicji~\ref{def:syntax-formula} są najmniejszymi
zbiorami zamkniętymi na wskazane konstrukcje.
Niech $T_P$ oznacza zbiór termów spełniających $P$.
Na mocy pierwszych dwóch założeń zawiera on wszystkie
zmienne i stałe, a na mocy ostatniego jest zamknięty
na tworzenie termów $f(\vec t)$. Minimalność zbioru
wszystkich termów daje $T_P=\operatorname{Term}(\mathcal L)$.
Dla formuł zbiór $F_Q$ formuł spełniających $Q$
zawiera wszystkie atomy i jest zamknięty na wszystkie
operacje z ich definicji; przez tę samą minimalność
jest więc zbiorem wszystkich formuł.
Tak otrzymujemy zasadę indukcji bez zakładania z góry
rozstrzygalności dowolnej własności $P$ lub $Q$.
\end{proof}

\begin{przyklad}[Liczba wolnych zmiennych jest skończona]
Każdy term i każda formuła zawierają skończenie wiele symboli,
zatem mają skończenie wiele wystąpień zmiennych.
Zbiór różnych zmiennych wolnych jest podzbiorem zbioru wszystkich
wystąpień zmiennych, a więc jest skończony. Wynika stąd, że
zawsze możemy wybrać \emph{świeżą} zmienną spośród
$v_0,v_1,\ldots$, niewystępującą w danym skończonym napisie.
\end{przyklad}